\documentclass[11pt]{article}
\usepackage{amsmath,amssymb,amsthm,booktabs,hyperref}
\usepackage[margin=1in]{geometry}
\newtheorem{theorem}{Theorem}
\title{A three-cycle and a four-period obstruction\protect\\for an equal-parameter delay recurrence}
\author{Principia Math In Progress}
\date{29 September 2026}
\begin{document}\maketitle
\begin{abstract}
We disprove MathDB 333489, the proposed period rule $p=k+3$ for
$z_{n+1}=\tau(1+z_{n-k})/(\tau-z_n)$.
The parameter $\tau=-2/3$ gives an exact three-cycle for every positive
delay $k\equiv1\pmod3$, so $k+3$ is not even a period of these solutions.
We also prove that for delay one and every nonzero complex parameter,
all solutions having period four already have period two. Thus the
conjecture also fails if interpreted as asserting the existence of an
orbit of least period $k+3$ for every delay.
\end{abstract}

\section*{Task 1: source scope and full counterexample}
Conjecture 3.2 on page 14 of Hassan's manuscript \cite{Hassan} proposes
$p=k+3$ when $\alpha=\beta=\gamma$ in
\[
 z_{n+1}=\frac{\alpha+\beta z_{n-k}}{\gamma-z_n}.
\]
Pages 1--2 allow arbitrary complex parameters and initial conditions;
Definition 3 distinguishes a period from its least (``prime'') period.
The conjecture has no extra stability or initial-value hypothesis.
We use nonconstant, everywhere-defined solutions, and address both
the universal-period and least-period existence readings.

\begin{theorem}\label{counter}
Let $k\ge1$ satisfy $k\equiv1\pmod3$ and set
$\alpha=\beta=\gamma=\tau=-2/3$. For every integer $n$, define
\[
 z_n=\begin{cases}
 0,&n\equiv0\pmod3,\\
 -1,&n\equiv1\pmod3,\\
 -2,&n\equiv2\pmod3.
 \end{cases}
\]
With the initial history $z_{-k},\ldots,z_0$ so defined, this is a
solution of the recurrence of least period three. The number $k+3$
is not a period of this solution.
\end{theorem}
\begin{proof}
Because $k\equiv1\pmod3$, the delayed entry is $z_{n-k}=z_{n-1}$.
All possible transitions are listed below.
\[
\begin{array}{c|r|r|c|r}
n\bmod3&z_{n-k}&z_n&\tau-z_n&
\tau(1+z_{n-k})/(\tau-z_n)\\\hline
0&-2&0&-2/3&-1\\
1&0&-1&1/3&-2\\
2&-1&-2&4/3&0
\end{array}
\]
No denominator is zero, and the last column is exactly $z_{n+1}$.
The three entries are distinct, so a shift is a period if and only if
it is divisible by three. But $k+3\equiv1\pmod3$.
In particular, $k=1$ gives an exact three-cycle instead of the proposed
period four, disproving the conjecture. Real parameters and initial
values are included in the source's complex domain.
\end{proof}

\section*{There is no orbit of least period four at delay one}
\begin{theorem}\label{four}
Fix any $\tau\in\mathbb C\setminus\{0\}$ and $k=1$. Every
everywhere-defined solution of period four of
\[
 z_{n+1}=\frac{\tau(1+z_{n-1})}{\tau-z_n}
\]
has period two. More precisely, its repeating entries are $a,b,a,b$
with $ab=-\tau$ and $a,b\ne\tau$.
Conversely, every such pair gives a solution of period dividing two.
\end{theorem}
\begin{proof}
Write a hypothetical four-periodic solution as $a,b,c,d,a,b,c,d,\ldots$.
Clearing its nonzero denominators gives
\begin{align*}
b(\tau-a)&=\tau(1+d),&
c(\tau-b)&=\tau(1+a),\\
d(\tau-c)&=\tau(1+b),&
a(\tau-d)&=\tau(1+c).
\end{align*}
Adding the first and third equations, and then the second and fourth,
gives $ab+cd=-2\tau=bc+ad$. Consequently
\[
 (a-c)(b-d)=0.
\]
After a cyclic relabeling if necessary, we may assume $a=c$.
The second equation shows that $a\ne0$, since $a=c=0$ would imply
$\tau=0$. Subtracting the second and fourth equations now gives
$a(d-b)=0$, hence $b=d$. The first equation then reduces to
$ab=-\tau$. Conversely, if $ab=-\tau$ and both denominators are
nonzero, then
\[
 b(\tau-a)=\tau(1+b),\qquad
 a(\tau-b)=\tau(1+a),
\]
which verifies the alternating solution.
\end{proof}

Theorem~\ref{four} applies also when $\tau$ is nonreal, including the
common parameter displayed in the source's numerical tables. It makes
no assertion about attraction, and does not rely on finite numerical
trajectories. Theorem~\ref{counter} already disproves the stated
universal period rule even when ``period'' is not required to be least.

\section*{Four proposed follow-up tasks}
\begin{enumerate}
\item[Task 2.] Classify the nonconstant three-cycles at delay one as
functions of the common complex parameter, including all forbidden poles.
\item[Task 3.] Determine for which delays and nonzero common parameters
an orbit of least period $k+3$ actually exists.
\item[Task 4.] Determine the stability of the periodic orbits in the
preceding classifications, separating existence from attraction.
\item[Task 5.] Reassess the source's numerical period tables using
specified initial histories, rigorous residual bounds, and certified
exclusion of shorter periods or long transients.
\end{enumerate}

\begin{thebibliography}{9}
\bibitem{Hassan} Sk.~Sarif Hassan,
\emph{Dynamics of the complex rational delay recursive sequence
$z_{n+1}=(\alpha+\beta z_{n-k})/(\gamma-z_n)$},
\href{https://arxiv.org/abs/1606.08887v1}{arXiv:1606.08887v1} (2016),
Definitions 1--3, pp.~1--2, and Conjecture~3.2, p.~14.
\end{thebibliography}
\end{document}
